\documentclass[11pt,a4paper,french]{article}

%---------------------------
% Packages
%---------------------------
\usepackage[utf8]{inputenc}
\usepackage[T1]{fontenc}
\usepackage[french]{babel}
\usepackage{amsmath, amssymb, amsfonts, amsthm}
\usepackage{mathrsfs}
\usepackage{enumitem}
\usepackage{hyperref}
\usepackage{geometry}
\usepackage{fancyhdr}
\usepackage{titlesec}
\usepackage{caption}
\usepackage{graphicx} % Insertion d'images
\usepackage{interval} 

%---------------------------
% Marges
%---------------------------
\geometry{top=2.5cm, bottom=2.5cm, left=1.5cm, right=1.5cm}

%---------------------------
% Style de page
%---------------------------
\pagestyle{fancy}
\fancyhead{}
\fancyhead[L]{\leftmark}
\fancyfoot[C]{\thepage}

%---------------------------
% Style des titres
%---------------------------
\titleformat{\section}{\normalfont\Large\bfseries}{\thesection}{1em}{}
\titleformat{\subsection}{\normalfont\bfseries}{\thesubsection}{1em}{}
\titleformat{\subsubsection}{\normalfont\small\bfseries}{\thesubsubsection}{1em}{}

%---------------------------
% Liens
%---------------------------
\hypersetup{
    colorlinks=false,
    linkcolor=red,
    urlcolor=blue,
    pdftitle={Documentation}, % À modifier
    pdfauthor={},
    pdfsubject={},
    pdfkeywords={},
}

%---------------------------
% Environnement
%---------------------------

\renewcommand{\b}[1]{\textbf{#1}}
\newcommand{\sautligne}{\vspace{0.5cm}}

\newcommand{\image}[2]{\begin{center} \includegraphics[width=#2cm]{./img/#1.png}\end{center}}

\newcommand{\test}{\b{question test}}

%---------------------------
% Début du document
%---------------------------
\begin{document}

\begin{titlepage}
    
    % --- Logo ---
    \begin{minipage}{0.20\textwidth}
        \image{logo_ensea}{7}
    \end{minipage}
    \vspace{4cm}

    \centering

    % --- Titre ---
    \hrule height 2pt
    \vspace{0.8cm}
    \begin{center}
        {\Huge \textbf{Forest Crawler - Documentation}}\\
        \vspace{0.3cm}
        {\Large } % À remplir
    \end{center}
    \vspace{0.8cm}
    \hrule height 2pt

    \vspace{2.5cm}

    % --- Informations ---
    {\large MABROUK Sofia, SAYAVONG Melvin} 

    \vspace{1cm}

    {\large ENSEA, 2025 - 2026, Semestre 6}

    \vfill
\end{titlepage}

\tableofcontents
\newpage

\section{Présentation}

\section{Class Main (peut-être continuer)}
\noindent La classe Main est la classe permettant d'exécuter tout le jeu, il s'occupe principalement de la gestion de la fenêtre.
\image{Pasted Image}{10}
Cette partie du code correspond aux paramètres de la fenêtre:
\begin{itemize}
    \item \verb|window.setDefaultCloseOperation(JFrame.EXIT_ON_CLOSE):| permet de déterminer le comportement du bouton de fermeture de la fenêtre.
    \item \verb|window.setResizable(false):| permet de savoir si la taille de la fenêtre est modifiable ou non.
    \item \verb|window.setTitle("Forest Crawler"):| correspond au nom de la fenêtre.
\end{itemize}

\section{Class Game Engine}
Cette classe s'occupe de la gestion du comportement des touches à l'aide d'un listener. Elle sert uniquement à configurer les touches pour le déplacement du joueur. 
\subsection{Variables}
Cette classe permet de gérer le comportement des touches du clavier à l'aide d'un listener. 
Dans notre cas, on utilise les touches du clavier uniquement pour se déplacer donc on créé 4 variables booléennes :
\begin{itemize}
    \item \verb|upPressed| 
    \item \verb|downPressed| 
    \item \verb|leftPressed| 
    \item \verb|rightPressed| 
\end{itemize}

\subsection{keyTyped}
\image{keyTyped}{7}
Cette méthode permet de savoir si une touche a été appuyé puis relâché. Elle ne sert pas pour ce projet, mais est 
automatiquement créée par IntelliJ dans la classe KeyListener (qui est implémentée par GameEngine) avec keyPressed et keyReleased.

\subsection{keyPressed}
\image{keyPressed}{7}
Cette méthode permet de savoir si les touches liées au déplacement ont été appuyé.

\subsection{keyReleased}
\image{keyReleased}{7}
Cette méthode permet de savoir si les touches liées au déplacement ont été relâché.

\section{RenderEngine}

Cette classe correspond à la configuration du rendu du jeu au sein de la fenêtre comme la taille des éléments du jeu.
Ce rendu dépend de deux groupes de paramètres de position: \verb|screenXXX|, une position relative à la fenêtre de jeu, et 
\verb|worldXXX|, une position par relative à la map du jeu, qui est 'out of bounds' de l'écran.
\subsection{Variables}

\subsubsection{Screen settings}
Voici la liste des différentes variables pour cette section:
\begin{itemize}
    \item \verb|originalTileSize| : Taille réelle des cases
    \item \verb|scale| : Échelle d'agrandissement
    \item \verb|tileSize| : Taille des cases lors de l'affichage
    \item \verb|maxScreenCol| : Nombre de colonnes affichées sur la fenêtre de jeu
    \item \verb|maxScreenRow| : Nombre de lignes affichées sur la fenêtre de jeu
    \item \verb|screenWidth| : Longueur de la fenêtre de jeu
    \item \verb|screenHeight| : Hauteur de la fenêtre de jeu
\end{itemize}

\subsubsection{World settings}
Voici la liste des différentes variables pour cette section:
\begin{itemize}
    \item \verb|maxWorldCol| : Nombre de colonnes de la map de jeu
    \item \verb|maxWorldRow| : Nombre de lignes de la map de jeu
    \item \verb|worldWidth| : Longueur maximum de la carte affichée
    \item \verb|worldHeight| : Hauteur maximum de la carte affichée
\end{itemize}

\subsubsection{FPS}
La variable \verb|FPS| correspond au nombre d'images par secondes.

\subsection{RenderEngine}
\image{RenderEngine}{15}
Cette méthode permet de définir les différents paramètres de la fenêtre.

\subsection{startGameThread}
\image{startGameThread}{10}
Cette méthode permet d'initialiser le thread pour le jeu.

\section{Class Player}
Cette classe permet de gérer tous les attributs et comportements liés au joueur comme sa position, son déplacement ou son sprite.

\subsection{Variables (description + à compléter)}
\begin{itemize}
    \item \verb|RenderEngine re|: instance de RenderEngine
    \item \verb|GameEngine keyH|: instance de GameEngine
    \item \verb|screenX|: position X du joueur sur la carte.
    \item \verb|screenY|: position Y du joueur sur la carte.
\end{itemize}

\subsection{Player}

\image{PlayerConstructor}{10}
Cette méthode correspond au constructeur de l'objet Player

\begin{itemize}
    \item \verb|solidArea| : Rectangle placé sur le joueur, qui sert de zone de collision (cf. CollisionCheck)
    \item \verb|screenX, screenY| : Positionnement du joueur au centre de la fenêtre de jeu
\end{itemize}

\subsection{setDefaultValues}
\image{setDefaultValues}{10}
Cette méthode permet de définir les valeurs par défaut lorsqu'on lance le jeu.

\subsection{getPlayerImage}
\image{getPlayerImage}{15}
Cette méthode permet de chercher les sprites du joueur en fonction de sa direction. Il y a 2 sprites par direction pour donner l'illusion du mouvement lors du déplacement.

\subsection{update (mettre image)}
Cette méthode permet de vérifier à tout instant les évènements de déplacement et de collision ainsi que de mettre à jour le sprite.

\subsection{draw (mettre une image)}
Cette méthode permet de mettre à jour la bonne image en fonction de la direction, en animant le mouvement.

\section{Class Sprite}
Cette classe permet de configurer le comportement des sprites.
\image{sprite}{12}
Voici la description des différentes variables:
\begin{itemize}
    \item \verb|worldX, worldY| : Position X et Y du sprite sur la carte
    \item \verb|speed| : Vitesse du sprite sur la carte
    \item \verb|up, down, right, left, idle| : Variables permettant d'utiliser les bonnes images en fonction de la direction
    \item \verb|direction| : Direction du sprite parmis [up, down, left, right]
    \item \verb|spriteCounter| : Compteur permettant d'actualiser les images pour l'animation
    \item \verb|spriteNumber| : Permet de basculer entre les différentes images lors d'une animation
    \item \verb|solidArea| : Hitbox du sprite
    \item \verb|collisionOn| : Activation de la collision
\end{itemize}

\section{Class CollisionCheck}
Cette classe permet de faire la gestion des collisions avec les différents sprite.

\subsection{CollisionCheck}
Constructeur pour définir le RenderEngine.

\subsection{checkTile (à faire)}

\section{Tile}
Cette classe permet de gérer l'affichage des cases ainsi que le comportement de celle-ci vis-à-vis du joueur. 
La class Tile contient 4 variables:
\begin{itemize}
    \item \verb|image| : prends l'image correspondante
    \item \verb|collision| : Permet de définir si l'image devra être une collision ou non
    \item \verb|isOverhead| : Permet de définir si l'image devra être au dessus du joueur ou non, (cf. méthode draw de TileManager)
    \item \verb|width, height| : Définit la taille de la case
\end{itemize}

\section{TileManager (à faire)}

\end{document}
